% \documentclass[a4paper,11pt,phdthesis,twoside,oneandhalfspace]{cssethesis}  %% <-- uncomment this string if you want to use old-school **dvips**
\documentclass[a4paper,11pt,phdthesis,twoside,oneandhalfspace,pdflatex]{cssethesis} %% <--- uncomment this string if you want to use PDFLaTeX

% * Include the option "pdflatex" above if you want to use pdflatex rather than dvips
% * Include the option "oneside" if you don't want formatting for two-sided  printing.
% * include ::
%                 = singlespace}{\AtBeginDocument{\singlespacing}}
%                 = oneandhalfspace}{\AtBeginDocument{\onehalfspacing}}
%                 = doublespace}{\AtBeginDocument{\doublespacing}}
% * Include the option "litreview" above if this is a literature review.
% * Include the option "nocoursecode" so that the numerical course code is suppressed after the course name.
% * Include option "thesisdraft" to get a timestamp and "Draft" message in   the footer
% * Include option "thesispsdraft" to get a timestamp and "Draft" message in
%   the footer, along with a grey "DRAFT" in the margin. Note: this only
%   works with latex, not pdflatex

%%%% LATEX STUFF %%%
\usepackage[T2A]{fontenc} 
\usepackage[british]{babel}

\usepackage{amssymb,amsmath,amsfonts,geometry,cite,enumerate,float}

\usepackage{setspace} %%% for additional spacing 
\onehalfspacing %%% <- Increase spacing for the proofread

\graphicspath{{images/}}  %% <-- where you store all your images. 
%%%% LATEX STUFF %%%




%%%% Thesis STUFF %%%
\thesisauthor{\Large John Doe}
\thesisauthorlastname{Doe}
% \thesisauthorpreviousdegrees{Dipl. Eng.} % Optional
\thesisdepartment{\bfseries School of Electrical Engineering \\ and Computer Science} % Optional.

\thesismonth{August} % Optional. Current month is used if this is not set
\thesisyear{2013} % Optional. Current year is used if this is not set

\thesistitle{On the development of a very important stuff}
\thesissupervisor{Prof. Woo Hoo}
%%%% Thesis STUFF %%%


%%%%%%%%%%%%%%%%%%%%%%%
\begin{document}
%%%%%%%%%%%%%%%%%%%%%%%


%%%%%%%%%%%%%%%%%%%%%%%%%%%%%%%%%%%%%%%%%%%%%%%%%%%%%%%%%%%%%%%%%%%%%%%%%%%%%%
%%
%% Front matter 
%%
\frontmatter					% start the thesis front matter.

\thesistitlepage				% Generate the title page.

\thesiscopyrightpage			% Generate the copyright page.

\thesisdeclarationpage			% generate the declaration page (optional).

\begin{thesisacknowledgments}	% generate the acknowledgements page (optional).

\lipsum[20-22]

\end{thesisacknowledgments}


\tableofcontents				% Generate a table of contents.

% \listoftables					% Generate a list of tables (optional).
% \listoffigures					% Generate a list of figures (optional).
%%%%%%%%%%%%%%%%%%%%%%%%%%%%%%%%%%%%%%%%%%%%%%%%%%%%%%%%%%%%%%%%%%%%%%%%%%%%%%


\begin{thesisabstract}

Very briefly summarise your work and what the thesis is all about. 


\end{thesisabstract}

%%%%%%%%%%%%%%%%%%  START::: List of Publications %%%%%%%%%%%%
\chapter{Publications and Outcomes}
The material  presented in this Thesis has been already published, or
accepted for publication, in peer-reviewed journals and conferences. Some of the
material has been recently submitted for publication. The list of publications is provided below.



\section*{Articles in peer-reviewed journals }

	\begin{description}
	\item[ A and B,]
	XXXXXXXXXX
     .  \emph{Journal on YYYYYYYYYY}


	\end{description}







\section*{Conference proceedings}

	\begin{description}
	\item[ A and B,]
	XXXXXXXXXX
     .  \emph{Journal on YYYYYYYYYY}


	\end{description}






%%%%%%%%%%%%%%%%%% ## END::: List of Publications %%%%%%%%%%%%




%%%%%%%%%%%%%%%%%%%%%%%%%%%%%%%%%%%%%%%%%%%%%%%%%%%%%%%%%%%%%%%%%%%%%%%%%%%%%%
%%
%% Main matter 
%%
\mainmatter	% start the thesis body.



%%%%%%%%%%%%%%%%%%%  START::: Introduction  %%%%%%%%%%%%
\chapter{Overview}


What the heck this thesis is about?! Give a very condensed summary.



\section{Review of previous works}
Unless you are an outright genius, your work must have been based on a work of your predecessors - cite them! Like this\cite{mita1977zeros}. Or, if there are many other works, like this\cite{astrom2011computer,annaswamy2001adaptive,doyle1979robustnesswithobservers}.


\section{Contributions of this Thesis}
Say what you have done (if anything) with your thesis - what is the actual contribution (problems posed and solved) with this thesis?


\section{Thesis Overview}
Provide a short overview of the contents of each chapter, chapter-by-chapter. Obviously, it is written when everything else is done.




\chapter{Introduction}

\section{Some good ideas about writing your thesis in LaTeX}

It goes without arguing that the Ph.D.~thesis in hard science
\emph{must} be written in LaTeX, and not in that dumbed-down
pointy-clicky behemoth MS Word/OpenOffice for office plankton.

The fact is: Word/Openoffice is useful for small, insignificant,
unstructured and unimportant documents. LaTeX, on the other hand,
is a powerful programming language for technical people and their
complex, well-structured, math-ridden and picture-rich documents.

The following tips and tricks can (and will!) save you countless
sleepless nights and oceans of tears:

\begin{itemize}
\item
  separate each chapter of your thesis in its own file, and then use
  \verb|\input{myfile}| command to link then in a master document
  \emph{thesisPhD.tex}
\item
  put pictures in a separate subdirectory/subfolder, preferable
  splitting them chapter-wise;
\item
  pick up either dvips or pdflatex,
  and don't mix them!
  The difference:
  \begin{itemize}
  \item
    use \textbf{dvips} if you have a lot of vector figures (from
    Gnuplot, for example) and vector graphics (in EPS format). You can
    always convert them to PDF using a command: \verb|epstopdf  foo.ps|
  \item
    use \textbf{pdflatex} in all other cases (if you don't need to
    convert your work to rtf/doc, or you don't know what EPS is).
  \end{itemize}
\item
  use \textbf{bibtex} or \textbf{natbib} for all your citations, it
  will save you countless buckets of tears
\item
  use the \verb|\label| command in LaTeX wisely:
  \begin{itemize}
  \item
    give names to different items with appropriate prefixes, e.g.,
    \verb|\label{eq:maxwell}| is a better reference to the Maxwell equations
    than \verb|\label{123}|
  \item
    use a colon (:) for references in chapters and sections, like
    \verb|\label{chap:intro:sec:review}|
  \end{itemize}
\item
  use version control system (Subversion/Mercurial/Git/whatever -
  they are FREE!) with LaTeX to:
  \begin{itemize}
  \item
    be able to roll back any change or recover a
    section/chapter/paragraph at your advisor's second thought (aka
    ``you know, a month ago we had a section here that it looked really nice\ldots{}'')
  \item
    to backup your whole thesis (especially true with distributed
    version control systems like Mercurial or Git) and sync your work
    between home/uni;
  \item
    to upload the code into a private repository on the Web (yet
    another backup).
  \end{itemize}
\item
  use bug/issues tracking capabilities in GitHub.com or bitbucket.com
  to flag the areas to discuss / improve.
\end{itemize}
\ldots{} and the rest of ideas from your thesis advisor.


\section{Mathematics}
Any \textit{good} thesis contains at least some mathematical expressions. 

In optimization we want to minimize or maximize an objective function over a set of decision variables subject to constraints. There are different ways to format optimization problems, but the preferable way is to follow an excellent book ``Convex Optimization'' by Stephen Boyd and Lieven Vandenberghe. A general optimization problem has the form:

\begin{equation*}
 \begin{aligned}
 & \underset{x}{\text{minimize}}
 & & f_0(x) \\
 & \text{subject to}
 & & f_i(x) \leq b_i, \; i = 1, \ldots, m.
 \end{aligned}
\end{equation*}

If we want to perform quadratic programming, the notation will be the following:
\begin{equation}\label{eq:qpoptimisation}
 \begin{aligned}
 & \underset{x}{\text{minimize}}
 & & J(x) = \frac{1}{2} \mathbf{x}^T Q \mathbf{x} + c^T \mathbf{x} \\
 & \text{subject to} 
 & & A\mathbf{x} \leq \mathbf b \mbox{  (inequality constraint)} \\
 & 
 & &     E\mathbf{x} = \mathbf d  \mbox{  (equality constraint)}
 \end{aligned}
\end{equation}

Eq.\ref{eq:qpoptimisation} is a rather sophisticated \LaTeX formulae. Unlike the tabular environment, in which you can specify the alignment of each column, in the aligned environment, each column (separated by \verb|&|) has a default alignment, which alternates between right and left-aligned. Therefore, all the odd columns are right-aligned and all the even columns are left-aligned.




\section{Tables}
An example of a simple table is provided below in Table~\ref{tab:example}. You may use \verb+booktabs+ package for more advanced/beautiful tables.

\begin{table}[ht]
	\begin{center}
	\begin{tabular}{lcr}
	23121 & 1212 & 232 \\ \hline  
	cat & frog & dog
	\end{tabular}
	\end{center}
\caption{An example of a table}
\label{tab:example}
\end{table}

See Table~\ref{tab:example}.


\section{Pictures}
An example picture is in Fig.~\ref{fig:examplepic}. You may also use floating figures with wrap-around (use the \verb|wrapfig| package, although they often create a lot of troubles.

\begin{figure}[ht!]
\centering\includegraphics[width=0.9\linewidth]{phd101310s}
\caption{An example picture.}
\label{fig:examplepic}
\end{figure}

Also you may consider a \verb|subfig| package if you want to put figures side-by-side and reference them separately (like Fig. 3a or Fig.3b).





\chapter{Conclusions and Future work}\label{chap:summary}



\appendix % all \chapter{..} commands after this will generate appendices

\chapter{This appendix should get a letter}
\label{app:example}
An appendix before the backmatter gets an automatically generated letter by
which it can be referred to. This is Appendix~\ref{app:example}.

\chapter{Simulation Source Code}
You may want to investigate the \texttt{lgrind} program and package if you
wish to include source code in your thesis

\backmatter % start the thesis back matter - the numbering of chapters stops here.



\chapter{Last Thing} % Appendices after the \backmatter command do not  get a letter
This sort of appendix has no letter. 



\pagestyle{plain}

\bibliographystyle{unsrt} 
\bibliography{biblio/my}  %%% this is a link to the Bibliography files
\end{document}